%\documentclass[DiscStateSpaceToolkit.tex]{subfiles} 
%\begin{document}
 
The Basic Real Business Cycle Model presented here simply extends the Stochastic Neo-Classical Growth Model to include an endogenous
choice of labour supply (of how much to work). This is an important difference in terms of the Toolkit we know have a '$d$' variable:
a choice variable that does not affect tomorrows state (in the Stochastic Neo-Classical Growth Model the only choice variable
was '$aprime$', next periods state). The following code solves the model. Below that are two further codes that 
providing examples of what one can then do with the solution of the value function iteration problem. The first demonstrates how to 
use the policy function created by the Value Function (Case 1) command to reproduce the relevant findings
of \cite*{AruobaFernandezVillaverdeRubioRamirez2006} looking at the sizes of numerical errors; this serves the added purpose of 
allowing you to easily experiment with how measures of the numerical errors change with the choice of grids. The second provide examples of how to 
use further Toolkit commands to simulate time series and calculate Standard Business Cycle Statistics.

The value function problem to be solved is,
\begin{equation*}
 V(k,z)=\sup_{k',l} \, \{ \frac{(c^{\theta} (1-l)^{1-\theta})^{\tau}}{1-\tau}+\beta*E[V(k',z')|z] \}
\end{equation*}
subject to
\begin{eqnarray}
 & c+i=exp(z) k^{\alpha} l^{1-\alpha} \\
 & k'=(1-\delta)k+i \\
 & z'=\rho z +\epsilon' \quad ,\epsilon \overset{iid}{\sim} N(0,\sigma_{\epsilon}^2)
\end{eqnarray}
where $k$ is physical capital, $i$ is investment, $c$ is consumption, $l$ is labour supply, $z$ is a productivity shock that follows an AR(1) process; 
$exp(z) k^{\alpha}$ is the production function, $\delta$ is the depreciation rate.

We will use the \href{appendix:TauchenMethod}{Tauchen Method} to discretize the AR(1) shock.

The codes implementing this model are available at  \href{https://github.com/vfitoolkit/VFItoolkit-matlab-examples}{github.com/vfitoolkit/VFItoolkit-matlab-examples}

The code $BasicRealBusinessCycleModel.m$ solves the model, using parallelization on GPU, and uses $BasicRealBusinessCycleModel\_{}ReturnFn.m$ which defines the return function. \\
There are also $BasicRealBusinessCycleModel\_{}BusinessCycleStatistics.m$ which reimplements solving the model and then goes on to use some of the other
Toolkit commands to simulate time series and calculate the Standard Business Cycle statistics. \\
Also $BasicRealBusinessCycleModel\_{}NumericalErrors.m$
which reimplements solving the model and then goes on to use the output to reproduces a number of relevant elements from 
Tables \& Figures in \cite*{AruobaFernandezVillaverdeRubioRamirez2006} relating to the accuracy of numerical solutions.
% 
% \textbf{Contents of StochasticNeoClassicalGrowthModel.m}
% \lstinputlisting[language=Matlab]{./Examples/MatlabCodes/BasicRealBusinessCycleModel.m}
% 
% \smallskip
% 
% \textbf{Contents of StochasticNeoClassicalGrowthModel$\_{}$ReturnFn.m}
% \lstinputlisting[language=Matlab]{./Examples/MatlabCodes/BasicRealBusinessCycleModel_ReturnFn.m}
% 
% \smallskip
% 
% \textbf{Contents of StochasticNeoClassicalGrowthModel$\_{}$BusinessCycleStatistics.m}
% \lstinputlisting[language=Matlab]{./Examples/MatlabCodes/BasicRealBusinessCycleModel_BusinessCycleStatistics.m}
% 
% \smallskip
% 
% \textbf{Contents of StochasticNeoClassicalGrowthModel$\_{}$NumericalErrors.m}
% \lstinputlisting[language=Matlab]{./Examples/MatlabCodes/BasicRealBusinessCycleModel_NumericalErrors.m}

% \textbf{Contents of AruobaFernandezVillaverdeRubioRamirez2006.m}
% \lstinputlisting[language=Matlab]{/home/kmarshallbanana/PhD/Matlab_Codes/DiscStateSpaceToolkit/AruobaFernandezVillaverdeRubioRamirez2006.m}
% 
% \smallskip
% 
% \textbf{Contents of AruobaFernandezVillaverdeRubioRamirez2006$\_{}$ReturnFn.m}
% \lstinputlisting[language=Matlab]{/home/kmarshallbanana/PhD/Matlab_Codes/DiscStateSpaceToolkit/AruobaFernadezVillaverdeRubioRamirez2006_ReturnFn.m}

%\end{document}
